\documentclass[11pt,a4paper,twocolumn]{article}

% === Core Packages ===
\usepackage[margin=2cm]{geometry}
\usepackage[utf8]{inputenc}
\usepackage[T1]{fontenc}
\usepackage{times}
\usepackage{amsmath,amssymb}
\usepackage{graphicx}
\usepackage{booktabs}
\usepackage{float}
\usepackage{enumitem}
\usepackage{tikz}
\usepackage{pgfplots}
\pgfplotsset{compat=1.18}
\usetikzlibrary{arrows.meta,shapes.geometric,positioning,calc,decorations.markings,fit,backgrounds}

% === Citation (IEEE style) ===
\usepackage[numbers,sort&compress]{natbib}
\bibliographystyle{IEEEtranN}

% === Hyperref ===
\usepackage[hidelinks]{hyperref}
\usepackage{xurl}

% === Settings ===
\setlength{\parindent}{1em}
\setlength{\parskip}{0.3em}

% === PDF metadata ===
\hypersetup{
  pdftitle={Ապակենտրոնացված հեղինակության ցանց՝ Հասարակության բնական կազմակերպման շրջանակ},
  pdfauthor={Pavel Kudrna},
  pdfsubject={Personix — a decentralized reputation network},
  pdfkeywords={Personix, decentralized reputation, reputation social network,
    decentralized identifiers, DID, meritocracy, censorship resistance},
}

% === Title ===
\title{\textbf{Ապակենտրոնացված հեղինակության ցանց՝\\Հասարակության բնական կազմակերպման շրջանակ}}
\author{Pavel Kudrna\\Personix Project --- \url{https://personix.org}\\Prague, Czechia\\Նախագիծ v1 --- Մարտ 2026}
\date{}

\begin{document}
\maketitle

% ============================================================
% ABSTRACT
% ============================================================
\begin{abstract}
Արդարության զգացումը՝ համոզմունքը, որ անարդարությունները ուղղվում են, իսկ արժանիքը՝ պարգևատրվում, հասարակության ամենաուժեղ միավորող ուժն է։ Երբ կառավարման կենտրոնացված ինստիտուտները չեն կարողանում ապահովել այս զգացումը, քանի որ իրավունքի պաշտպանության արժեքը գերազանցում է հենց իրավունքի արժեքը, սոցիալական համախմբվածությունը քայքայվում է։ Ներկայիս ինստիտուցիոնալ կառույցները չեն կարողանում լուծել վեճերի զգալի դասը հենց այնպես, ինչպես կենտրոնացված պլանավորումը ձախողվում է ռեսուրսների բաշխման հարցում՝ տեղեկատվության անհամաչափության, հետադարձ կապի բացակայության և որոշում կայացնողների՝ իրենց որոշումների հետևանքներից անջատվածության պատճառով։ Այս հոդվածը ներկայացնում է Հեղինակության սոցիալական ցանցը (RSN)՝ ապակենտրոնացված, անկաշառ, գրաքննության դիմացկուն հաղորդակցման շերտ, որը կառուցված է Ապակենտրոնացված նույնացուցիչների (DID) վրա և թույլ է տալիս կեղծանվանական մասնակիցներին հրապարակել, ստուգել և սպառել հեղինակության վերաբերյալ տեղեկատվություն իրական աշխարհի վարքագծի մասին։ Հիմնական մեխանիզմը հենվում է երեք նախագծման սկզբունքի վրա. (1)~անհամաչափ ծախսային կառուցվածք, որտեղ հեղինակության տվյալների ընթերցումն էժան է, բայց գրառումը պահանջում է ստուգելի ջանք. (2)~ստուգողի ընտրության ոչ դետերմինիստական արձանագրություն՝ հիմնված հետևողական հեշավորման վրա, որը գնագոյացնում է արմատական պնդումները սրընթաց աճող կրկնության ծախսերի միջոցով. և (3)~շուկայի կողմից պայմանավորված հեղինակության իշխանության մոդել, որտեղ մասնավոր ստուգողները գրավ են դնում իրենց կուտակած հեղինակությունը այն պնդումների ճշգրտության վրա, որոնք նրանք հաստատում են։ Ստացված ճարտարապետությունը ստեղծում է առաջացող մերիտոկրատիա առանց կենտրոնական համակարգման և հնարավոր է դարձնում աստիճանական, հետշրջելի անցումային ուղի գոյություն ունեցող պետական վարչարարական համակարգերից։
\end{abstract}

% ============================================================
% 1. INTRODUCTION
% ============================================================
\section{Ներածություն}

\subsection{Կենտրոնացված վստահության խնդիրը}

Ժամանակակից կառավարումը հենվում է հիմնարար ենթադրության վրա. որ կենտրոնացված ինստիտուտները կարող են օտար մարդկանց միջև վստահությունը միջնորդել մեծ մասշտաբով։ Դատարանները լուծում են վեճերը։ Ռեեստրերը հաստատում են սեփականությունը։ Կարգավորող մարմինները կիրառում են չափանիշներ։ Այս ինստիտուտները նախագծվել են դարաշրջանում, երբ տեղեկատվությունը սակավ էր, հաղորդակցությունը՝ դանդաղ, և իշխանության պատվիրակումը կենտրոնական մարմնին միակ իրագործելի համակարգման մեխանիզմն էր։ Կենտրոնացված կառավարումը, սակայն, ձախողվում է նույն կառուցվածքային պատճառներով, ինչ կենտրոնացված պլանավորումը՝ տեղեկատվության անհամաչափություն, հետադարձ կապի բացակայություն և որոշում կայացնողների՝ իրենց որոշումների հետևանքներից անջատվածություն։

Ենթադրությունը ձախողվում է, օրինակ, երբ ինստիտուցիոնալ միջնորդության արժեքը գերազանցում է վեճի արժեքը։ Դիտարկենք վարձակալ, ով պարզում է, որ իր վարձակալած բնակարանը չունի օրենքով պահանջվող էլեկտրական անվտանգության վկայական՝ կյանքին անմիջական վտանգ ներկայացնող պայման։ Վարձակալի իրավական իրավունքը՝ դուրս գալ պայմանագրից, միանշանակ է։ Սակայն այդ իրավունքը դատական համակարգի միջոցով պաշտպանելու արժեքը գերազանցում է գրավի և պարտք եղած վնասների դրամական արժեքը, նույնիսկ ներառելով հնարավոր օրենսդրական փոխհատուցումը~\cite{czcourtfees}։ Բանական պատասխանը կորուստը կլանելն ու դուրս գալն է։

Սա պաթոլոգիական բացառիկ դեպք չէ։ Դա վեճերի զգալի դասի կանխադրված փորձն է, որտեղ վնասը իրական է, բայց դրամական խաղադրույքն ընկնում է այն շեմից ցածր, որից բարձր ինստիտուցիոնալ պաշտպանությունը դառնում է տնտեսապես բանական։ Արդյունքը այնպիսի համակարգ է, որը կառուցվածքային առումով նպաստում է չար դերակատարներին. նրանք, ովքեր վնասում են ուրիշներին այս շեմից ցածր ծավալներով, կարող են գործել անպատիժ, քանի որ արդարություն փնտրելու արժեքը ընկնում է տուժած կողմի վրա։

Վերը նշված օրինակը վերցված է հեղինակի իրավական միջավայրից՝ չեխական քաղաքացիական իրավունքի համակարգից։ Այլ իրավական ավանդույթներում՝ նախադեպի վրա հիմնված ընդհանուր իրավունք, սովորութային իրավունք, խառը համակարգեր, արդարությունը ձախողվում է տարբեր ձևերով և տարբեր աստիճաններով՝ կախված իրավասության մշակութային և ընթացակարգային առանձնահատկություններից։ Ընթերցողները հավանաբար կճանաչեն համարժեք օրինակներ իրենց սեփական համատեքստից։ Ինչ որ կառուցվածքային առումով համընդհանուր է, դա օրինաչափությունն է. վեճեր, որոնց արժեքն ընկնում է տնտեսապես բանական պաշտպանության շեմից ցածր, ավարտվում են նրանով, որ կորուստը կլանում է տուժած կողմը։ RSN-ը կատարյալ արդարություն չի խոստանում. այն միայն նվազեցնում է այն կառուցվածքային առավելությունը, որը վայելում են չար դերակատարները, երբ ինստիտուցիոնալ պաշտպանությունը ոչ տնտեսական է։

\subsection{Ինչու գոյություն ունեցող լուծումները թերի են}

\textbf{Ապակենտրոնացված նույնականացման (DID) շրջանակները}~\cite{did2022} տրամադրում են ծածկագրական մեխանիզմներ ինքնիշխան նույնականության կառավարման համար, բայց հիմնականում կենտրոնանում են ինքնության ստուգման վրա՝ առանց անդրադառնալու վարքագծային հեղինակության ավելի լայն հարցին։

\textbf{Հոգեկապ նշանները (SBT)}~\cite{weyl2022} ընդգրկում են այն հասկացողությունը, որ հեղինակությունը չփոխանցելի է, բայց հենվում են բլոկչեյն ենթակառուցվածքի վրա՝ մասշտաբայնության սահմանափակումներով, և չեն անդրադառնում տեղեկատվության մուտքագրումը կառավարող տնտեսական խթաններին։ Հարակից հասկացությունները, ինչպիսիք են արժանիքի նշանները (օրինակ՝ Liberland), կիսում են նմանատիպ փիլիսոփայություն, բայց մնում են կապված կոնկրետ իրավասական շրջանակների հետ։

\textbf{Ապակենտրոնացված ինքնավար կազմակերպությունները (DAO)}~\cite{buterin2014} իրականացնում են կառավարումը խելացի պայմանագրերի միջոցով, բայց վերարտադրում են պլուտոկրատական դինամիկա, որտեղ ազդեցությունը համաչափ է կապիտալին։ Քառակուսի քվեարկությունը~\cite{buterin2019} մասամբ մեղմացնում է դա, բայց սահմանափակում է գործնական ընդունումը։ DAO-ները նաև բախվում են հիմնարար \emph{օրակուլի խնդրին}. խելացի պայմանագրերը չեն կարող հուսալիորեն ստուգել իրական աշխարհի վիճակները առանց վստահելի արտաքին աղբյուրի, և այս խնդիրը չունի ընդհանուր լուծում բլոկչեյն ճարտարապետության շրջանակներում։

Ոչ մի գոյություն ունեցող համակարգ չի համատեղում ապակենտրոնացումը, գրաքննության դիմադրությունը, կեղծանվանականությունը, մուտքի ստուգելի արժեքը և առաջացող համաձայնությունը։

\subsection{Ներդրումներ}

Այս հոդվածը կատարում է հետևյալ ներդրումները.
\begin{enumerate}[nosep]
\item \textbf{Հեղինակության սոցիալական ցանցի (RSN)} սահմանում՝ ապակենտրոնացված արձանագրություն, որն ընդլայնում է DID շրջանակը՝ երկկողմ հեղինակության փոխանակմանն աջակցելու համար։
\item \textbf{Ստուգողի ընտրության ոչ դետերմինիստական արձանագրություն}՝ հիմնված հետևողական հեշավորման վրա~\cite{karger1997}։
\item \textbf{Թանկ արմատականության սկզբունքը}, որը գնագոյացնում է պնդումներն ըստ ցանցի համաձայնությունից նրանց հեռավորության՝ երեք կուտակվող ծախսային ալիքների միջոցով։
\item \textbf{Հեղինակավոր իշխանության մոդել}՝ անհամաչափ խթաններով, որտեղ կեղծ հաստատումն արժենում է աղետալիորեն ավելի, քան օրինական վճարները։
\item \textbf{Աստիճանական անցումային ուղի}՝ զուգահեռ ֆիսկալ ենթակառուցվածքի միջոցով։
\end{enumerate}

% ============================================================
% 2. DESIGN PRINCIPLES
% ============================================================
\section{Նախագծման սկզբունքներ}

RSN ճարտարապետությունը սահմանափակված է հինգ չբանակցելի նախագծման սկզբունքով։

\textbf{Շարունակականություն և էվոլյուցիա.} Համակարգը գոյակցում է գոյություն ունեցող ինստիտուտների հետ անորոշ երկարության անցումային ժամանակահատվածում։ Անցումը պետք է լինի հետշրջելի յուրաքանչյուր փուլում։

\textbf{Կամավորություն և պատասխանատվություն.} Մասնակցությունը կամավոր է, բայց կամավորությունը զուգորդվում է պատասխանատվության հետ. մասնակիցը, ով ստանձնում է վերահսկողություն ինստիտուցիոնալ գործառույթի նկատմամբ, միաժամանակ ստանձնում է ուղեկցող պարտավորությունները։

\textbf{Բազմազանություն և մշակութային չեզոքություն.} Համաձայնությունն առաջանում է երկկողմ փոխազդեցություններից, ոչ թե համակարգի նախագծողների կողմից պարտադրված նախապես սահմանված կանոններից։

\textbf{Դիմացկունություն և գրաքննության դիմադրություն.} Ցանցը գրաքննության ենթարկելու արժեքը պետք է գերազանցի այն հանդուրժելու արժեքը։ Միայն լիակատար տոտալիտարիզմը՝ կործանարար քաղաքական և տնտեսական ծախսերով, կարող է ճնշել համակարգը։

\textbf{Համաձայնություն և կիրառում.} Համակարգը ներառում է մարդկային հակումները դեպի մանրախնդրություն, չարություն և վրիժառու բավարարվածություն որպես կրող հատկանիշներ։ Իրավախախտումը արգելված չէ. այն գնագոյացվում է։

% ============================================================
% 3. SYSTEM OVERVIEW
% ============================================================
\section{Համակարգի ընդհանուր նկարագիր}

\subsection{Հեղինակության սոցիալական ցանց}

RSN-ը ապակենտրոնացված, հասակ առ հասակ հաղորդակցման շերտ է, որտեղ մասնակիցները փոխանակում են ստուգելի տեղեկատվություն իրական աշխարհի վարքագծի մասին։ Դրա որոշիչ հատկությունները հետևյալն են՝ ապակենտրոնացված և չգրաքննվող ենթակառուցվածք. կեղծանվանական մասնակցություն DID-ների միջոցով. անհամաչափ ծախսային կառուցվածք (ընթերցումն էժան է, գրառումը՝ թանկ). ծածկագրական ստուգելիություն. և \textbf{ստանդարտների աջակցություն}՝ մեքենայաընթեռնելի ձևաչափեր ցանցով տարածվող տեղեկատվության համար, իշխանությունների առաջարկած ծառայությունների համար (պնդման կազմում, հաստատում, վկայի ծառայություն), և քաղաքականության ձևաչափի համար, ներառյալ հակառակ կողմի հեղինակությունը գնահատելու և քաղաքականության անհամապատասխանության (հայտարարված և իրական վարքագծի միջև) ռիսկը գնահատելու կանոնները։

\subsection{Ճարտարապետական բաղադրիչներ}

RSN-ը բաղկացած է չորս հիմնական բաղադրիչից (Նկ.~\ref{fig:architecture}).
\begin{enumerate}[nosep]
\item \textbf{DID շերտ.} Մասնակիցների ինքնությունները՝ հայտարարված կանոններով, հեղինակության մետատվյալներով և ցանցային երթուղավորմամբ։
\item \textbf{Պնդման արձանագրություն.} Կառուցվածքային ձևաչափ իրական աշխարհի իրադարձությունների մասին պնդումներ հրապարակելու համար։
\item \textbf{Ստուգման ցանց.} Ապակենտրոնացված արձանագրություն ստուգողներ նշանակելու համար՝ հետևողական հեշավորման միջոցով։
\item \textbf{Հեղինակության ագրեգացման շերտ.} Ծառայություններ, որոնք արտադրում են մարդկաընթեռնելի հեղինակության ամփոփումներ։
\end{enumerate}

\begin{figure}[ht]
\centering
\begin{tikzpicture}[
    layer/.style={draw, minimum height=0.7cm, align=center, font=\scriptsize, fill=black!8},
    arr/.style={<->, thick, gray!60}
]
\node[layer, minimum width=5cm] (agg) at (0,2.8) {\textbf{Ագրեգացում}\\Մեկնաբ.\ $\cdot$ Ռիսկ};
\node[layer, minimum width=2.2cm] (claim) at (-1.55,1.4) {\textbf{Պնդումներ}\\Կազմել};
\node[layer, minimum width=2.2cm] (verif) at (1.55,1.4) {\textbf{Ստուգում}\\Ընտրություն};
\node[layer, minimum width=5cm] (did) at (0,0) {\textbf{DID շերտ}\\Ինքնություն $\cdot$ Կանոններ $\cdot$ Միավորներ};

\draw[arr] (agg.south west) -- (claim.north);
\draw[arr] (agg.south east) -- (verif.north);
\draw[arr] (claim.east) -- (verif.west);
\draw[arr] (claim.south) -- (did.north west);
\draw[arr] (verif.south) -- (did.north east);
\end{tikzpicture}
\caption{RSN չորս շերտանոց ճարտարապետություն։}
\label{fig:architecture}
\end{figure}

\subsection{Մասնակիցների դերեր}

Ստուգման գործարքը ներառում է մինչև վեց տարբեր DID դեր (Նկ.~\ref{fig:roles})՝ \textbf{Թողարկող} ($DID_I$, ստեղծում և ներկայացնում է պնդումը), \textbf{Սուբյեկտ} ($DID_S$, այն DID-ը, որի մասին պնդումն է՝ կարող է համընկնել Թողարկողի հետ), \textbf{Իշխանություն} ($DID_A$, հաստատում է պնդումը վճարի դիմաց. կարող է նաև առաջարկել վկայի ծառայություններ՝ \emph{ընտրովի}), \textbf{Վկա} ($DID_{W_{1..n}}$, ստուգողի ազնվության ինկոգնիտո վերահսկիչ մարտահրավերի կոդերի միջոցով՝ \emph{ընտրովի}), \textbf{Ստուգող} ($DID_V$, ալգորիթմորեն ընտրված պնդումը հրապարակելու համար), և \textbf{RSN} (ցանցը, որտեղ հրապարակվում են ստուգված պնդումները)։ Ցանկացած դեր կարող է պատվիրակվել մեկ այլ DID-ի (Բաժին~7.4)։ Իշխանությունը սովորաբար միավորում է հաստատումը վկայի ծառայությունների հետ, բայց երկու գործառույթները տրամաբանորեն անկախ են։

\begin{figure}[ht]
\centering
\begin{tikzpicture}[
    role/.style={draw, rounded corners, minimum width=1.1cm, minimum height=0.45cm, align=center, font=\scriptsize, fill=black!8},
    opt/.style={draw, rounded corners, minimum width=1.1cm, minimum height=0.45cm, align=center, font=\scriptsize, fill=black!8, dashed},
    arr/.style={-{Stealth[length=3pt]}, thick},
    oarr/.style={-{Stealth[length=3pt]}, thick, dashed, gray},
    lbl/.style={font=\tiny, fill=white, inner sep=1pt}
]
\node[role] (I) at (0,2.4) {\textbf{Թողարկող}};
\node[role] (S) at (-2.5,2.4) {\textbf{Սուբյեկտ}};
\node[opt] (A) at (2.5,1.2) {\textbf{Իշխանություն}};
\node[opt] (W) at (-2.5,0) {\textbf{Վկա}};
\node[role] (V) at (2.5,-0.6) {\textbf{Ստուգող}};
\node[role, fill=black!5, draw=gray] (N) at (0,-1.8) {RSN};

\draw[arr] (I) -- node[lbl] {մասին} (S);
\draw[oarr] (I) -- node[lbl,sloped] {հաստատում} (A);
\draw[oarr] (I) -- node[lbl,sloped] {մարտահրավեր} (W);
\draw[arr] (I) -- node[lbl,sloped] {հարցում} (V);
\draw[oarr] (W) -- node[lbl] {վերահսկում} (V);
\draw[arr] (V) -- node[lbl] {հրապարակում} (N);
\end{tikzpicture}
\caption{DID դերերը ստուգման գործարքում։ Կետագիծ = ընտրովի (Իշխանություն, Վկա)։}
\label{fig:roles}
\end{figure}

\subsection{Անկաշառություն. չորս ուժ}

RSN-ի անկաշառությունը չի բխում մեկ մեխանիզմից, այլ չորս միաժամանակյա ուժից։ Ոչ մեկն ինքնին բավարար չէ. միայն դրանց համակցությունն է կաշառակերության փորձը դարձնում տնտեսապես անբանական։

\textbf{Անկանխատեսելի թիրախ.} Յուրաքանչյուր պնդման ստուգողն ընտրվում է ոչ դետերմինիստորեն՝ հետևողական հեշավորման միջոցով (Բաժին~5.3)։ Հարձակվողը չի կարող նախապես թիրախավորել կոնկրետ ստուգողի կաշառելու համար, քանի որ ստուգողի ինքնությունը պնդման բովանդակության և նախորդ կրկնությունների ֆունկցիա է, ոչ թե թողարկողի ընտրության։

\textbf{Հեղինակության լծակ՝ դանդաղ վեր, արագ վար.} Հեղինակությունը կարելի է ձեռք բերել միայն աստիճանաբար՝ կայուն հետևողական վարքագծի միջոցով։ Սակայն այն կարող է աղետալիորեն կորչել մեկ խարդախ հաստատմամբ (Բաժին~6.2)։ Այս անհամաչափությունը նշանակում է, որ տարիների ընթացքում կուտակված հեղինակությունը կարող է ոչնչանալ մեկ անազնիվ հաստատմամբ. օպտիմալ ռազմավարությունը մնում է կասկածելի պնդումները մերժելը։

\textbf{Հրապարակային խոստում.} Յուրաքանչյուր DID Կրող հրապարակայնորեն հայտարարում է իր քաղաքականությունը՝ մեքենայաընթեռնելի կանոնների մի շարք, որը պարտավորվում է հետևել։ Հայտարարության և իրական վարքագծի միջև շեղումը հայտնաբերելի է և գնագոյացվում է որպես հեղինակության հանցագործություն, ավելի ծանր, քան հիմքում ընկած խախտումը (Բաժին~7.3)։ Հետևաբար կեղծավորությունն ավելի թանկ արժե, քան ուղղակի անազնվությունը։

\textbf{Ցանցային հարձակման-ծախսի անհամաչափություն.} Ցանցը գրաքննության ենթարկելու կամ ապակայունացնելու արժեքն աճում է ոչ գծայնորեն՝ ցանցի չափի և խտության հետ, մինչդեռ այն հանդուրժելու արժեքը մնում է հաստատուն։ Որպեսզի գրաքննությունը մարի, կենտրոնացված դերակատարը ստիպված կլիներ այն կիրառել ամբողջ ցանցում՝ պահանջելով միջոցներ, որոնք անհամաչափ են ցանկացած պատկերացնելի օգուտի (Բաժին~10)։

% ============================================================
% 4. DECENTRALIZED IDENTITY MODEL
% ============================================================
\section{Ապակենտրոնացված նույնականության մոդել}

\subsection{DID հատուկ հատկություններով}

RSN-ը ընդլայնում է W3C DID Core բնութագիրը~\cite{did2022} հետևյալով՝ \textbf{Հայտարարված կանոններ} (մեքենայաընթեռնելի վարքագծային պարտավորություններ), \textbf{Հեղինակության մետատվյալներ} (ագրեգացված վարքագծային միավոր), և \textbf{Ցանցային երթուղավորում} (փոփոխական սոխաձև դարպաս, առանձին կայուն օղակի դիրքից)։

\subsection{Պնդման կենսապարբերաշրջան}

Պնդումն անցնում է վեց փուլով (Նկ.~\ref{fig:lifecycle})՝ կազմում, ընտրովի իշխանության հաստատում, ստուգողի ընտրություն հետևողական հեշավորման միջոցով, ստուգման որոշում (ընդունում կամ մերժում կրկնությամբ), հրապարակում, և հեղինակության թարմացում բոլոր կողմերի համար։

\begin{figure}[ht]
\centering
\begin{tikzpicture}[
    stp/.style={draw, rounded corners=2pt, minimum width=1.3cm, minimum height=0.45cm, align=center, font=\tiny, fill=black!5},
    arr/.style={-{Stealth[length=3pt]}, thick},
    lbl/.style={font=\tiny, fill=white, inner sep=1pt}
]
\node[stp] (comp) at (0,0) {Կազմել\\պնդում};
\node[stp, dashed] (auth) at (2.0,0) {Իշխանության\\հաստատում};
\node[stp] (hash) at (4.0,0) {Հեշ և\\օղակի քարտեզ};
\node[stp] (ver) at (2.0,-1.6) {Ստուգողի\\ստուգում};
\node[stp, double] (pub) at (0,-1.6) {Հրապարակված};
\node[stp] (rej) at (4.0,-1.6) {Մերժված\\ծախս $\uparrow$};

\draw[arr] (comp) -- (auth);
\draw[arr] (auth) -- (hash);
\draw[arr] (hash) -- (ver);
\draw[arr] (ver) -- node[lbl] {\checkmark} (pub);
\draw[arr] (ver) -- node[lbl] {$\times$} (rej);
\draw[arr] (rej) -- ++(0,0.8) -| (hash);
\end{tikzpicture}
\caption{Պնդման կենսապարբերաշրջանը կրկնության օղակով։}
\label{fig:lifecycle}
\end{figure}

\subsection{Կապը W3C DID Core-ի հետ}

RSN-ի նույնականության մոդելը W3C DID Core բնութագրի ճիշտ վերաբազմություն է~\cite{did2022}։ Բոլոր RSN DID-ները վավեր W3C DID-ներ են։ RSN-ին հատուկ ընդլայնումները կոդավորված են որպես DID փաստաթղթի հատկություններ, որոնք սահմանված են հատուկ DID մեթոդի բնութագրում՝ համատեղելի գոյություն ունեցող ենթակառուցվածքի հետ, ինչպիսիք են ION~\cite{buchner2021} և Ceramic~\cite{ceramic2023}։

% ============================================================
% 5. VERIFICATION AND PROPAGATION
% ============================================================
\section{Տեղեկատվության ստուգում և տարածում}

\subsection{Տեղեկատվության մուտքագրման արժեքը}

RSN-ի հիմնական տնտեսական սկզբունքը ընթերցման և գրառման միջև անհամաչափությունն է։ Հեղինակության տվյալների ընթերցումը մոտենում է զրոյական սահմանային ծախսի այն մասնակիցների համար, ովքեր DID ցանցի մասն են և ունեն իրենց հեղինակությանը համապատասխան հասանելիություն. նորեկները պետք է աստիճանաբար վաստակեն ավելի լայն հասանելիություն (տես հակամակաբուծություն)։ Գրառումը՝ հասկացված որպես հեղինակության մշտական նյութականացում ցանցում, ոչ թե որպես միանվագ տեխնիկական ակտ, պահանջում է ստուգելի կուտակային ներդրում երեք չափումների վրա՝ \textbf{ժամանակ} (կյանքի տարիներ ծախսված հետևողական վարքագծի, կատարված պարտավորությունների և մշակված հարաբերությունների վրա), \textbf{էներգիա} (ազնիվ վարվելակերպի, գործընկերությունների նկատմամբ հոգածության և երկարաժամկետ վստահելիության վրա ներդրված ջանք), և \textbf{փող} (սեփական անվան մեջ ներդրում՝ մասնագիտական զարգացում, սեփական աշխատանքի որակ, պատճառված վնասների փոխհատուցում)։ Հեղինակությունը հնարավոր չէ ակնթարթորեն գնել. այն ցմահ կուտակման արդյունք է, որը համակարգը պարզապես դարձնում է տեսանելի և փոխանցելի համատեքստերի միջև։ Արձանագրության միանվագ տեխնիկական ծախսերը (ծածկագրական ստորագրություններ, ստուգողի վճարներ) աննշան են այս հիմքում ընկած ներդրման համեմատ։

\subsection{Սոցիալական գրաֆ և կապի խորություն}

RSN-ը հիմնականում սոցիալական ցանց է. յուրաքանչյուր DID ավելացնում է կոնտակտներ (այլ DID-ներ), որոնք թույլ են տալիս կապը։ Այս կոնտակտներն ունեն իրենց սեփական կոնտակտները, և այդպես շարունակ։ Ալգորիթմական ստուգողի որոնումը գործում է կապակցված կոնտակտների կարգավորելի $h$ խորության շրջանակներում (օրինակ՝ $h=3$՝ սեփական ուղղակի կոնտակտները, նրանց կոնտակտները և կոնտակտների կոնտակտները)։ Այս ճարտարապետությունը չի պահանջում մեկ գլոբալ բլոկչեյն. այն բնականորեն նպաստում է համայնքների/կլաստերների առաջացմանը՝ այլ համայնքների հետ համընկնումներով։ Յուրաքանչյուր DID մասնակից պետք է ունենա հրապարակված \textbf{քաղաքականություն}՝ մեքենայաընթեռնելի կանոնների մի շարք, որը կառավարում է, թե ինչպես են գնահատվում մուտքային հարցումները։ Քաղաքականության խախտումները գրանցվում են ցանցում որպես բացասական արտեֆակտներ։

\subsection{Ալգորիթմական ստուգողի ընտրություն}

Ստուգման արձանագրությունն օգտագործում է հետևողական հեշավորում~\cite{karger1997} (Նկ.~\ref{fig:verifier})։ Ստուգումը վճարովի ծառայություն է. ստուգողները գործում են վճարի դիմաց՝ ստեղծելով շուկա, որտեղ մրցակցությունը պայմանավորում է գնագոյացումը, արագությունը և հուսալիությունը։

\textbf{Քայլ 1.} Պնդման փաստաթուղթը միակցվում է նախորդ կրկնության հեշավորման արդյունքի հետ (կամ $\varnothing$՝ առաջին կրկնության դեպքում) և հեշավորվում է.
\begin{equation}
\text{pos} = f_h(\text{claim} \| \text{prev\_hash})
\label{eq:hash}
\end{equation}

Ալգորիթմը պետք է ներառի բոլոր նախորդ հարցումների և պատասխանների \emph{ամբողջական ուղին}՝ ոչ միայն նախորդ կրկնության հեշը։ Յուրաքանչյուր ստուգողի ամբողջական պատասխանը (ընդունում, մերժում, առաջարկվող գին, պատճառ) կցվում է աճող փաստաթղթին՝ ստուգելի ծածկագրական ստորագրությունների միջոցով յուրաքանչյուր քայլում։

\textbf{Քայլ 2.} Հեշը քարտեզագրվում է հետևողական հեշավորման օղակի $N$ դիրքերի, հավասարաչափ բաշխված (օրինակ՝ $N=4$-ի համար՝ $+0^{\circ}, +90^{\circ}, +180^{\circ}, +270^{\circ}$)։ Յուրաքանչյուր դիրքից ժամասլաքի ուղղությամբ որոնումն ընտրում է ամենամոտ DID-ը որպես ստուգողի թեկնածու (Նկ.~\ref{fig:multipoint})։ $N$-ը փոփոխական պարամետր է, որը կարող է կախված լինել ընթացիկ ակտիվ DID-ների տոկոսից, օրվա ժամից կամ ցանցի պատմական արձագանքունակությունից։

\textbf{Քայլ 3.} Ստուգողն ընդունում է (հրապարակում է՝ գրավ դնելով հեղինակությունը) կամ մերժում (վերադառնալով Քայլ~1-ին մերժման հեշով)։ Երբ հանգույցը չի արձագանքում՝ չնայած առցանց կարգավիճակ հայտարարելուն, հաջորդ հարցումը պետք է ներառի բոլոր $N$ նախորդ ստուգողի թեկնածուների բոլոր պատասխանները։

\textbf{Հակամակաբուծություն.} Թիրախ DID-ները կարող են սահմանել վճարներ կամ հասանելիության սահմանափակումներ քիչ հեղինակություն ունեցող նոր DID-ների հարցումների համար։ Այս աստիճանական մեխանիզմը (ոչ բինար) ստիպում է նորեկներին կառուցել հեղինակություն իսկական գործունեության միջոցով՝ նախքան ուրիշների տվյալներն ազատ սպառելը։

\begin{figure}[ht]
\centering
\begin{tikzpicture}[
    box/.style={draw, rounded corners=2pt, minimum width=1.3cm, minimum height=0.45cm, align=center, font=\scriptsize, fill=black!5},
    dec/.style={draw, diamond, aspect=2.2, minimum width=0.9cm, align=center, font=\scriptsize, fill=black!5},
    arr/.style={-{Stealth[length=3pt]}, thick},
    lbl/.style={font=\tiny, fill=white, inner sep=1pt}
]
\node[box] (iss) at (0,0) {Թողարկող};
\node[box] (hash) at (0,-1.1) {$f_h$(claim $\|$\\prev\_hash)};
\node[box] (ring) at (0,-2.2) {Օղակի քարտեզ\\ժամասլաքով};
\node[box, dashed] (spam) at (0,-3.2) {Հակասպամ\\ավանդ};
\node[dec] (check) at (0,-4.4) {Քաղաքականության\\ստուգում};
\node[box] (rej) at (2,-5.6) {Հաջորդ\\կրկնություն};
\node[box] (fee) at (0,-5.6) {Վերջնական վճար =\\$f$(տվյալ, հեղ., քաղ.)};
\node[box, double] (pub) at (0,-6.8) {Հրապարակված};

\draw[arr] (iss) -- (hash);
\draw[arr] (hash) -- (ring);
\draw[arr] (ring) -- (spam);
\draw[arr] (spam) -- (check);
\draw[arr] (check) -- node[lbl] {$\times$} (rej);
\draw[arr] (check) -- node[lbl] {\checkmark} (fee);
\draw[arr] (fee) -- (pub);
\draw[arr] (rej.east) -- ++(0.5,0) |- (hash.east);
\end{tikzpicture}
\caption{Ստուգողի ընտրության արձանագրություն։ Հակասպամ ավանդն ընտրովի է և կարող է վերադարձվող լինել։ Վերջնական վճարը վճարվում է միայն ընդունման դեպքում։}
\label{fig:verifier}
\end{figure}

Յուրաքանչյուր մերժում կցում է ստուգողի ամբողջական պատասխանը (ներառյալ պատճառներն ու պայմանները) պնդման փաստաթղթին՝ ընդլայնելով դրա բովանդակությունը։ Ընդլայնված փաստաթղթից ոչ դետերմինիստորեն հաշվարկվում է նոր հեշ՝ քարտեզագրվելով օղակի տարբեր դիրքի և ընտրելով տարբեր ստուգող։ Ամբողջական ուղու փաստաթղթավորումը պարտադիր է. թողարկողը չի կարող կանխատեսել կամ ազդել հաջորդ ստուգողի վրա։ Եթե ստուգողը անտեսում է հարցումը՝ չնայած համատեղելի հայտարարված քաղաքականությանը, վկաները (եթե առկա են) գրանցում են դա, և թողարկողը կարող է կցել վկայի ապացույցը վերջնական հրապարակման ժամանակ։

\begin{figure}[ht]
\centering
\begin{tikzpicture}[scale=0.65]
% Ring
\draw[thick] (0,0) circle (2.2cm);
% DID nodes on ring
\foreach \a in {10,55,80,120,150,195,230,260,305,340} {
  \fill[black!40] (\a:2.2) circle (1.5pt);
}
% Hash offset arrow pointing to initial position
\draw[-{Stealth[length=3pt]}, thick, black!60] (0,0) -- node[font=\tiny,above,sloped,pos=0.6] {$f_h$} (37:2.2);
\fill[black] (37:2.2) circle (2pt);
\node[font=\tiny,anchor=south west] at (37:2.35) {$h_0$};
% N=4 points offset from h0 by +0, +90, +180, +270
\foreach \offset/\l in {0/P1,90/P2,180/P3,270/P4} {
  \pgfmathsetmacro{\ang}{37+\offset}
  \node[fill=black, circle, inner sep=1.5pt] (p\offset) at (\ang:2.2) {};
  \node[font=\tiny,font=\bfseries] at (\ang:2.75) {\l};
  % clockwise search arc
  \draw[-{Stealth[length=2.5pt]}, thick, gray] (\ang:2.2) arc (\ang:\ang+30:2.2);
}
\node[font=\scriptsize, align=center] at (0,0) {Հեշ\\Օղակ};
\node[font=\tiny, align=center] at (0,-3.2) {$h_0 = f_h(\text{claim})$, ապա $+0^{\circ},+90^{\circ},+180^{\circ},+270^{\circ}$\\Յուրաքանչյուր կետ $\rightarrow$ ժամասլաքով որոնում};
\end{tikzpicture}
\caption{Բազմակետ ընտրություն ($N{=}4$)։ Հեշ $h_0$-ը սահմանում է շեղումը. 4 կետ հավասարաչափ բաշխված $h_0$-ից։}
\label{fig:multipoint}
\end{figure}

\subsection{Վկայի արձանագրություն}

\textbf{Վկա} դերը ապահովում է ինկոգնիտո հաշվետվողականություն ստուգողի ազնվության համար՝ ծածկագրական մարտահրավեր-կոդի արձանագրության միջոցով.

\textbf{Քայլ W1.} Հարցումատուն ստորագրում է ստուգման հարցումը և ուղարկում այն $N_w$ վկաների՝ հարցումատուի սոցիալական ցանցից կոնտակտներ կամ վճարի դիմաց գործող մասնագիտացված վկայի ծառայություն մատուցողներ։

\textbf{Քայլ W2.} Յուրաքանչյուր վկա $w_i$ ստորագրում է ամբողջ ստորագրված հարցումը, պահպանում է սկզբնական ստորագրությունը $\sigma_i$, և հաշվարկում մարտահրավերի կոդ $c_i = h(\sigma_i)$։ Մարտահրավերի կոդը կցվում է հարցմանը։

\textbf{Քայլ W3.} Հարցումը, այժմ կրելով $N_w$ մարտահրավերի կոդ, ուղարկվում է ստուգողին։ Ստուգողը դիտում է կոդերը, բայց \emph{չի կարող որոշել}, թե ովքեր են վկաները կամ արդյոք կոդերը համապատասխանում են իրական ստորագրություններին։

\textbf{Քայլ W4 (ազնիվ ստուգող).} Եթե ստուգողն արձագանքում է իր հայտարարված քաղաքականությանը համապատասխան, մարտահրավերի կոդերը մնում են անթափանց։ Վկաները երբեք չեն բացահայտվում։ Լրացուցիչ տվյալներ չեն հրապարակվում։

\textbf{Քայլ W5 (քաղաքականության խախտում).} Եթե ստուգողը խախտում է իր քաղաքականությունը (անտեսում է հարցումը, արձագանքում է անհետևողականորեն), հարցումատուն պահում է վկաների սկզբնական ստորագրությունները $\sigma_i$։ Դրանք կարող են հրապարակվել որպես ուղեկցող տվյալ. ցանկացած ոք կարող է ստուգել $c_i = h(\sigma_i)$՝ ապացուցելով, որ վկան ներկա էր։ Հարցումատուն կարող է հրապարակել ինքնուրույն. ստուգողն արդեն ստորագրել է մարտահրավերի կոդերը իր պատասխանում։

\textbf{Բլֆի հատկություն.} Հարցումատուն կարող է որոշ կամ բոլոր մարտահրավերի կոդերի փոխարեն պատահական արժեքներ դնել։ Ստուգողը չի կարող տարբերել իսկական կոդերը (հենված բարձր հեղինակություն ունեցող իշխանությունների վրա) աղմուկից։ Այս \emph{անորոշությունն} է կիրառման մեխանիզմը. յուրաքանչյուր հարցում կրում է այն ռիսկը, որ հարգված իշխանությունը ինկոգնիտո հետևում է։ Այս ճնշումը պահպանելու արժեքը գրեթե զրո է (պատահական թվեր գեներացնելը), մինչդեռ ստուգողի համար անազնվության հնարավոր արժեքն աղետալի է։ Ազնիվ վարքագիծը խրախուսվում է նույնիսկ իրական վկաների բացակայության դեպքում։

\textbf{Իշխանությունը որպես վկա.} Իշխանությունը, որը հաստատում է պնդում, կարող է միավորել վկայի ծառայությունները. «Ես հաստատում եմ ձեր պնդումը և ծառայում որպես ինկոգնիտո վկա ստուգման ընթացքում»։ Սա բնական բիզնես մոդել է. Իշխանությունն արդեն ունի համատեքստը։ Այնուամենայնիվ, երկու դերերը տրամաբանորեն անկախ են։

\subsection{Թանկ արմատականության սկզբունքը}

Երեք ծախսային ալիք միաժամանակ կուտակվում են (Նկ.~\ref{fig:radicalism}).

\textbf{Ալիք 1. Կրկնության ծախս.} Յուրաքանչյուր մերժում պարտադրում է նոր կրկնություն՝ սպառելով ժամանակ և ռեսուրսներ։ Գծային աճ։

\textbf{Ալիք 2. Փաստաթղթի ուռճացում.} Յուրաքանչյուր կրկնություն ավելացնում է ստուգողի ամբողջական պատասխանը պնդման փաստաթղթին։ Աճը գծային է, բայց հիմքի շեղումով. սկզբնական բեռը (պնդման բովանդակություն, իշխանության հաստատում, ծածկագրական ստորագրություններ) արդեն ունի ոչ չնչին չափ $B_0$։ Ընդհանուր չափը $k$ կրկնությունից հետո՝ $S(k) = B_0 + k \cdot \Delta$, որտեղ $\Delta$-ն միջին պատասխանի չափն է։

\textbf{Ալիք 3. Ստուգողի ռիսկի հավելավճար.} Ռիսկը սուբյեկտիվ է։ Ստուգողը գնահատում է, թե արդյոք հարցում ներկայացնող DID-ի հայտարարված քաղաքականությունները հակասում են ստուգողի սեփական առևտրային, սոցիալական կամ քաղաքական շահերին։ Հավելավճարը կարող է աճել էքսպոնենցիալ՝ ստուգողի «իդեալից» ընկալվող հեռավորության հետ. $P(d) = \alpha \cdot e^{\beta d}$, որտեղ $d$-ն ստուգողի սուբյեկտիվ հեռավորության չափանիշն է։ Անձնական արգելված սահմանից այն կողմ ստուգողը կարող է ամբողջովին հրաժարվել։

\begin{figure}[ht]
\centering
\begin{tikzpicture}
\begin{axis}[
    width=\columnwidth,
    height=4.5cm,
    xlabel={\footnotesize Արմատականություն},
    ylabel={\footnotesize Հարաբերական ծախս (log)},
    ymode=log,
    xmin=0, xmax=100,
    ymin=1, ymax=10000,
    grid=major,
    grid style={gray!20},
    legend style={at={(0.03,0.97)},anchor=north west,font=\tiny,draw=gray!50},
    tick label style={font=\tiny},
    label style={font=\footnotesize},
]
\addplot[black, thick, domain=0:100, samples=50] {1 + 0.3*x};
\addlegendentry{Կրկնության ծախս (գծային)}
\addplot[black, thick, dashed, domain=0:100, samples=50] {1 + 0.15*x};
\addlegendentry{Փաստաթղթի ուռճացում (գծային)}
\addplot[black, thick, dotted, line width=1.2pt, domain=0:100, samples=100] {exp(0.08*x)};
\addlegendentry{Ռիսկի հավելավճար (էքսպոնենցիալ)}
\end{axis}
\end{tikzpicture}
\caption{Ծախսի սրընթաց աճ երեք ալիքներով։ Ռիսկի հավելավճարը գերիշխում է բարձր արմատականության դեպքում։}
\label{fig:radicalism}
\end{figure}

Կարևոր է, որ սա գրաքննություն չէ։ Ոչ մի պնդում արգելված չէ։ Համակարգը գնագոյացնում է ռիսկը (Աղյուսակ~\ref{tab:costs})։

\begin{table}[ht]
\centering
\caption{Ծախսերի համեմատություն ըստ պնդման տեսակների։}
\label{tab:costs}
\footnotesize
\begin{tabular}{@{}lcccc@{}}
\toprule
\textbf{Տեսակ} & \textbf{Կրկն.} & \textbf{Փաստ.} & \textbf{Վճար} & \textbf{Ընդամ.} \\
\midrule
Վստահելի & $k_{\min}$ & $B_0$ & $P_{\min}$ & Ցածր \\
Վիճելի & $k_{\text{med}}$ & $B_0{+}k\Delta$ & $\alpha e^{\beta d}$ & Չափավոր \\
Արմատական & $k_{\max}$ & $\gg B_0$ & $\gg P_{\min}$ & Շատ բարձր \\
Խարդախ & $\to\infty$ & --- & --- & Չհրապարակելի \\
\bottomrule
\end{tabular}
\end{table}

% ============================================================
% 6. REPUTATION AND RISK
% ============================================================
\section{Հեղինակություն և ռիսկի գնահատում}

\subsection{Հեղինակության կուտակման մոդել}

Հեղինակությունն ունի երեք հատկություն՝ \textbf{դանդաղ կուտակում} (աստիճանական աճ հետևողական վարքագծի միջոցով), \textbf{արագ ոչնչացում} (մեկ խարդախությունը կարող է աղետալիորեն նվազեցնել միավորը), և \textbf{անհատական գնահատում} (յուրաքանչյուր սպառող կողմ կարող է կիրառել իր սեփական ագրեգացման ֆունկցիան)։

\subsection{Հեղինակավոր իշխանություններ}

Հեղինակավոր իշխանությունը վերանայում է պնդումները և, եթե բավարարված է, համաստորագրում դրանք՝ գրավ դնելով իր սեփական հեղինակությունը (Նկ.~\ref{fig:authority})։ Անհամաչափությունը միտումնավոր է. հեղինակություն ձեռք բերելը դանդաղ է (աստիճանական $+\delta$ յուրաքանչյուր ազնիվ հաստատման դիմաց), մինչդեռ դրա կորուստը աղետալի է ($-\Delta \gg \delta$ յուրաքանչյուր կեղծ հաստատման դիմաց. պատկերավոր օրինակ՝ $+2\%$-ի դիմաց $-70\%$)։ Օպտիմալ ռազմավարությունը միշտ կասկածելի պնդումները մերժելն է։ $\delta$-ի և $\Delta$-ի կոնկրետ արժեքները համակարգի պարամետրեր են։

\begin{figure}[ht]
\centering
\begin{tikzpicture}[
    box/.style={draw, rounded corners=2pt, minimum width=2cm, minimum height=0.6cm, align=center, font=\scriptsize, fill=black!5},
    arr/.style={-{Stealth[length=4pt]}, thick}
]
\node[box] (exam) at (0,0) {Իշխանությունը քննում է\\Հեղինակություն՝ $R$};
\node[box] (true) at (-2,-1.3) {Ճշմարիտ՝ $R{\to}R{+}\delta$\\Ավելի շատ հաճախորդ};
\node[box] (false) at (2,-1.3) {Կեղծ՝ $R{\to}R{-}\Delta$\\Հաճախորդները հեռանում};
\node[box] (insight) at (0,-2.6) {Ձեռքբերում՝ դանդաղ ($+\delta$)\\Կորուստ՝ ակնթարթ ($-\Delta$)};

\draw[arr] (exam) -- node[left,font=\tiny] {ճշմարտ.} (true);
\draw[arr] (exam) -- node[right,font=\tiny] {սուտ} (false);
\draw[arr] (true) -- (insight);
\draw[arr] (false) -- (insight);
\end{tikzpicture}
\caption{Հեղինակավոր իշխանություն՝ անհամաչափ խթաններ։}
\label{fig:authority}
\end{figure}

\subsection{Բազմաչափ հեղինակություն}

Հեղինակությունը սկալյար չէ, այլ վեկտոր՝ բազմաթիվ չափումների վրա՝ ֆինանսական հուսալիություն, անձնական բարեխղճություն, մասնագիտական իրավասություն, քաղաքացիական ներգրավվածություն և այլն (Նկ.~\ref{fig:reputation-radar})։ Տարբեր գնահատման մատակարարներ նախագծում են այս վեկտորը տարբեր կերպ՝ կախված համատեքստից. վարկատուն առաջնահերթ է համարում ֆինանսական հուսալիությունը, գործատուն՝ մասնագիտական իրավասությունը, հարևանը՝ քաղաքացիական վարքագիծը։ Չկա մեկ «ճիշտ» հեղինակության միավոր. կան միայն տեսանկյուններ։

\begin{figure}[ht]
\centering
\begin{tikzpicture}[scale=0.55]
\foreach \a/\l in {90/Ֆինանսական,162/Բարեխղճություն,234/Մասնագիտական,306/Քաղաքացիական,18/Սոցիալական} {
  \draw[gray!40] (0,0) -- (\a:2.5);
  \node[font=\tiny, align=center] at (\a:3.1) {\l};
  \foreach \r in {0.8,1.6,2.4} {
    \fill[black!15] (\a:\r) circle (0.5pt);
  }
}
\draw[thick, fill=black!10] (90:2.0) -- (162:1.2) -- (234:1.8) -- (306:0.9) -- (18:1.5) -- cycle;
\draw[thick, dashed] (90:1.4) -- (162:2.1) -- (234:0.8) -- (306:2.0) -- (18:1.1) -- cycle;
\node[font=\tiny] at (0,-3.3) {Հոծ՝ DID A \quad Կետագիծ՝ DID B};
\end{tikzpicture}
\caption{Բազմաչափ հեղինակության վեկտորներ։ Տարբեր DID-ներ ցուցաբերում են տարբեր պրոֆիլներ չափումների վրա։}
\label{fig:reputation-radar}
\end{figure}

\subsection{Ժողովրդավարացված ռիսկի գնահատում}

RSN-ը ժողովրդավարացնում է համակարգված ռիսկի գնահատումը՝ ունակություն, որ ներկայումս կենտրոնացված է ֆինանսական հաստատություններում, բայց կիրառելի է բանկային գործից շատ ավելի հեռու։ Արդյունաբերական կոնգլոմերատներ, որոնք գնահատում են մատակարարների հուսալիությունը, ապահովագրական ընկերություններ, որոնք գնագոյացնում են վարքագծային ռիսկը, միաձուլումների և ձեռքբերումների պատշաճ ջանասիրություն, արտադրության մեջ սարքավորումների ծառայության ժամկետի գնահատում. ամենուր, որտեղ հակառակ կողմի ռիսկը պահանջում է գնահատում, RSN-ը տրամադրում է ապակենտրոնացված, շուկայի կողմից պայմանավորված այլընտրանք։ Մեկնաբանման ծառայությունների շուկան թարգմանում է հումք բազմաչափ հեղինակության տվյալները գործնական ամփոփումների՝ դարձնելով հակառակ կողմի գնահատումը հասանելի ցանկացած մասնակցի սահմանային ծախսով։

% ============================================================
% 7. CONSENSUS AND DISPUTE RESOLUTION
% ============================================================
\section{Համաձայնություն և վեճերի լուծում}

\subsection{Ապակենտրոնացված արդարության մոդել}

RSN-ը վերաձևակերպում է արդարությունը որպես երկկողմ բանակցության խնդիր։ Մշտական, ստուգելի հեղինակության վնասի սպառնալիքն ավելի արդյունավետ զսպող միջոց է, քան դատական վարույթները, որոնք տուժած կողմը չի կարող իրեն թույլ տալ։

\subsection{Առաջացող համաձայնություն}

Յուրաքանչյուր մասնակից պահպանում է անձնական բավարարվածության շեմ։ Ցանցի ագրեգացված համաձայնությունն առաջանում է որպես հազարավոր երկկողմ բանակցությունների վիճակագրական միջին (Նկ.~\ref{fig:consensus})։ Համաձայնությունից շատ բարձր (բարբարոսական) պատասխանը թանկ է հրապարակել։ Համաձայնությանը մոտ (համաչափ) պատասխանն էժան է։ Համաձայնությունը կանոն չէ. այն գնային ազդանշան է։

\begin{figure}[ht]
\centering
\begin{tikzpicture}[
    person/.style={draw, rounded corners=2pt, minimum width=1.5cm, minimum height=0.5cm, align=center, font=\scriptsize, fill=black!5},
    arr/.style={-{Stealth[length=4pt]}, thick}
]
\node[person] (a) at (-2,0) {Խիստ\\(բարձր)};
\node[person] (b) at (0,0) {Պրագմատիկ\\(միջին)};
\node[person] (c) at (2,0) {Ներողամիտ\\(ցածր)};
\node[person, fill=black!10, minimum width=3cm] (cons) at (0,-1.2) {Ցանցի համաձայնություն\\= վիճ.\ միջին};
\node[person] (exp) at (-2,-2.4) {Բարբարոսական\\թանկ};
\node[person, double] (prop) at (0,-2.4) {Համաչափ\\էժան};
\node[person] (neg) at (2,-2.4) {Անփույթ\\թանկ};

\draw[arr] (a) -- (cons);
\draw[arr] (b) -- (cons);
\draw[arr] (c) -- (cons);
\draw[arr] (cons) -- (exp);
\draw[arr] (cons) -- (prop);
\draw[arr] (cons) -- (neg);
\end{tikzpicture}
\caption{Առաջացող համաձայնություն անհատական բավարարվածության շեմերից։}
\label{fig:consensus}
\end{figure}

\subsection{Պատժի չափորոշում}

Յուրաքանչյուր DID Կրող հրապարակայնորեն հայտարարում է պատասխանները իրավախախտումների կոնկրետ կատեգորիաների համար։ Անհամաչափ խիստ կամ մեղմ հայտարարությունները ներգրավում են բացասական հեղինակության ազդանշաններ։ Համաչափ հայտարարություններն ընդունվում են առանց շփման՝ ինքնաչափորոշվող պատժի համակարգ։

Համաձայնության հայտարարություններն իրենք ենթակա են կեղծավորության հայտնաբերման. համաձայնության դիրքին հայտարարաբար հավատարիմ մնալը՝ միևնույն ժամանակ անհետևողականորեն վարվելով, հեղինակության հանցագործություն է, ավելի ծանր, քան հիմքում ընկած շեղումը, քանի որ դա դրսևորում է կանխամտածված խաբեություն։

\subsection{Պատվիրակում}

DID-ի ներսում բոլոր գործունեությունները՝ բացառությամբ մեկի, կարող են պատվիրակվել մեկ այլ DID-ի։ \textbf{Սուբյեկտի դերը՝ այն DID-ը, որի վարքագծի մասին է պնդումը, չի կարող պատվիրակվել. այն անփոխանցելիորեն կապված է ինքնության կրողի հետ, ում վերաբերում է պնդումը։} Մյուս ակտիվ դերերը՝ Թողարկող, Իշխանություն, Վկա, Ստուգող, կարող են փոխանցվել նշանակված պատվիրակ DID-ի, լինի դա ստուգման, պնդման ներկայացման, համաձայնության մասնակցության կամ վեճերի լուծման համար։ Պատվիրակումը հետկանչելի է ցանկացած պահի։ Սա հնարավոր է դարձնում մասնագիտացումը (օրինակ՝ մասնագիտական ստուգման ծառայություններ), մինչդեռ Կրողը պահպանում է վերջնական վերահսկողությունն ու պատասխանատվությունը։ Պատվիրակումը կիրառվում է ամբողջ համակարգում և էական է մասշտաբայնության համար։

\subsection{Անհատական բարոյականությունից դեպի առաջացող հասարակական պայմանագիր}

Յուրաքանչյուր DID-ի հայտարարված քաղաքականությունը ներկայացնում է անհատական բարոյականության ձևակերպում. ինչ է Կրողը համարում ընդունելի, ինչպես է արձագանքում կոնկրետ վարքագծերին, ինչ է պահանջում հակառակ կողմերից։ Այս քաղաքականությունները պարտադրված չեն համակարգի նախագծողի կողմից. դրանք ընտրվում են յուրաքանչյուր մասնակցի կողմից և արտահայտվում մեքենայաընթեռնելի ձևով։

Այս ձևակերպումը հնարավոր է դարձնում եռաստիճան առաջացումը։ Առաջին, անհատական բարոյականությունը կոդավորվում է որպես \textbf{քաղաքականություն}՝ հայտարարված կանոնների, շեմերի և պատասխանի ֆունկցիաների մի շարք, որին DID-ը պարտավորվում է հետևել։ Երկրորդ, ցանցում բոլոր քաղաքականությունների ագրեգատն արտադրում է \textbf{առաջացող էթիկա}՝ վիճակագրորեն դիտարկելի նորմեր, որոնք ոչ մի առանձին մասնակից չի նախագծել, բայց որ առաջանում են հազարավոր անհատական բարոյական դիրքորոշումների փոխազդեցությունից~\cite{durkheim1893}։ Երրորդ, այս առաջացող էթիկան, արտահայտված որպես երկկողմ գնային ազդանշանների միջոցով կիրառվող ընդհանուր վարքագծային նորմեր, կազմում է \textbf{չափելի հասարակական պայմանագիր}՝ ոչ թե տեսական կառույց, որը ոչ ոք չի ստորագրել~\cite{rawls1971}, այլ էմպիրիկորեն դիտարկելի կանոնների մի շարք, որը հենված է իրական վարքագծի վրա։

Այս գործընթացն արտացոլում է բարոյական հիմքերի տեսության եզրակացությունները~\cite{haidt2012}. մարդկային բարոյականությունը ինտուիտիվ է, բազմակարծ և մշակութային առումով փոփոխական։ RSN-ը չի պահանջում բարոյական համաձայնություն որպես ներդրում. այն արտադրում է վարքագծային համաձայնություն որպես արդյունք։ Ստացված հասարակական պայմանագիրը ստատիկ չէ. այն զարգանում է, երբ մասնակիցները միանում են, հեռանում և հարմարեցնում իրենց քաղաքականությունները՝ ի պատասխան ցանցի հետադարձ կապի։ Ի տարբերություն դասական հասարակական պայմանագրի տեսության, այս պայմանագիրը հետկանչելի է, դիտարկելի և կիրառելի առանց ինքնիշխանի։

% ============================================================
% 8. ECONOMIC MODEL
% ============================================================
\section{Տնտեսական մոդել}

Նախորդ բաժինները նկարագրեցին գործիք՝ RSN-ի ստուգման մեխանիզմները, ծախսային կառուցվածքը և առաջացող համաձայնությունը։ Այս բաժինը նկարագրում է այս գործիքը ֆիսկալ և կառավարման անցմանը կիրառելու մեթոդաբանություն։ Գործիքն ընդհանուր նշանակության է. ստորև ներկայացված մեթոդաբանությունը մեկ հնարավոր կիրառություն է։

\subsection{Խթանների կառուցվածք}

Հեղինակությունը որպես կապիտալ ստեղծում է ուղղակի խթաններ ազնիվ վարքագծի համար։ Նորմալ գործունեության ընթացքում մասնակիցները բնականորեն կառուցում են հեղինակություն ամենօրյա գործարքների և կատարված պարտավորությունների միջոցով։ Հիմնական ծախսն ուղղված է \emph{բացասական} պնդումներին՝ ուրիշների պատճառած վնասի գրանցմանը։ Այս անհամաչափությունը խրախուսում է օգտակար, հուսալի վարքագիծ. ազնիվ լինելը լրացուցիչ ոչինչ չարժե, մինչդեռ վնասակար լինելը ստեղծում է փաստաթղթավորված հետք, որը ուրիշները կվճարեն պահպանելու համար։ Կեղծավորությունը՝ կանոններ հայտարարելը՝ առանց դրանց հետևելու, ամենաթանկ ձախողման ձևն է, քանի որ դրսևորում է կանխամտածված խաբեություն։

\subsection{Զուգահեռ ֆիսկալ համակարգ}

Երեք բաղադրիչ հնարավոր են դարձնում մրցակցային ճնշումը. (1)~\textbf{Պարզ հարկ}՝ միասնական համաչափ դրույքաչափ առանց բացառությունների, սկզբում մի փոքր ցածր գոյություն ունեցող փաստացի դրույքաչափից. (2)~\textbf{Ծախսերի էլեկտրոնային գրանցում (ESR)}՝ չեխական EET մոդելի հակադարձում, որպեսզի քաղաքացիները վերահսկեն պետական ծախսերը. (3)~\textbf{Քաղաքացու ուղղորդած հարկերի բաշխում}՝ սկսած մի քանի տոկոսից առաջին տարում և հասնելով, մեկ տասնամյակ անց, ցածր երկնիշ թվերից մինչև ամբողջական միգրացիա դեպի քաղաքացու ուղղորդած համակարգ՝ կախված ընդունման տեմպից։ Իրական տեմպը կախված է պետական ծառայությունների ազատ-շուկայական այլընտրանքների առաջացման արագությունից և պետության՝ անցումին համագործակցելու պատրաստակամությունից. ժամանակի ֆունկցիան կարիք չունի գծային լինելու, և չկա որևէ պատճառ վերին սահման դնելու այն բանի, թե որտեղ կարող է ի վերջո հասնել ընդունումը։

% ============================================================
% 9. MIGRATION PATH
% ============================================================
\section{Անցումային ուղի}

Անցումն ընթանում է երեք փուլով (Նկ.~\ref{fig:migration})՝ \textbf{Փուլ~1. Ծածկույթ} (RSN-ը որպես տեղեկատվական շերտ, պետությունը միակ մատակարարն է), \textbf{Փուլ~2. Մրցակցություն} (RSN-ը տրամադրում է գործառական այլընտրանքներ, երկակի համակարգի օգտագործում), և \textbf{Փուլ~3. Անցում} (RSN-ը գերազանցում է պետական համարժեքներին, կամավոր միգրացիա)։ Անցումը հետշրջելի է յուրաքանչյուր փուլում։

\begin{figure}[ht]
\centering
\begin{tikzpicture}[
    phase/.style={draw, rounded corners=3pt, minimum width=1.8cm, minimum height=0.8cm, align=center, font=\scriptsize, fill=black!5},
    arr/.style={-{Stealth[length=5pt]}, thick}
]
\node[phase] (p1) at (0,0) {\textbf{Փուլ 1}\\Ծածկույթ};
\node[phase] (p2) at (2.2,0) {\textbf{Փուլ 2}\\Մրցակցություն};
\node[phase] (p3) at (4.4,0) {\textbf{Փուլ 3}\\Անցում};

\draw[arr] (p1) -- (p2);
\draw[arr] (p2) -- (p3);
\draw[arr, dashed, gray] (p3.south) -- ++(0,-0.6) -| node[below,font=\tiny,pos=0.25] {հետդարձ} (p2.south);
\end{tikzpicture}
\caption{Եռափուլ միգրացիա՝ հետշրջելի յուրաքանչյուր փուլում։}
\label{fig:migration}
\end{figure}

Ճնշման հիմնական մեխանիզմը ֆիսկալ է. երբ պայմանական հարկատուները հասնում են «տնտեսապես նշանակալի փոքրամասնության $X$»՝ խումբ, որը բավական մեծ է, որ դրա դեմ ճնշումն առաջացնում է ոչ չնչին տնտեսական վնաս, և քաղաքացիական անհնազանդությունը դառնում է վստահելի սպառնալիք, պետությունը մտնում է բանակցության մեջ։ Ցուցադրության համար մենք օգտագործում ենք $X \approx 15\%$ ՀՆԱ-ից, բայց փաստացի շեմը կախված է կոնկրետ տնտեսությունից։

% ============================================================
% 10. SECURITY ANALYSIS
% ============================================================
\section{Անվտանգության վերլուծություն}

Մենք դիտարկում ենք հակառակորդի հինգ կատեգորիա՝ անհատ չար դերակատարներ, կազմակերպված խմբեր, կորպորատիվ հակառակորդներ, պետական մակարդակի հակառակորդներ և արձանագրության մակարդակի հակառակորդներ։

\textbf{Սիբիլի դիմադրություն}~\cite{douceur2002}։ Հեղինակություն կառուցելը պահանջում է կայուն, ստուգելի փոխազդեցություն։ Հաջողված Սիբիլ հարձակման արժեքն աճում է գծայնորեն կեղծ ինքնությունների հետ և քառակուսայնորեն թիրախային հեղինակության մակարդակի հետ։ Զուգահեռ բազմաթիվ DID-ներ շահագործելը հնարավոր է, բայց թանկ՝ նախագծմամբ. յուրաքանչյուր ինքնություն պետք է ինքնուրույն կուտակի հեղինակություն իսկական գործունեության միջոցով։ Ազատ հասարակություններում այս արժեքը զսպում է պատահական չարաշահումը. բռնապետություններում զուգահեռ DID-ները դառնում են գոյատևման գործիք՝ հնարավոր դարձնելով բաժանարար դիմադրություն, ընդհատակյա համակարգում և ավելի անվտանգ սև-շուկայական նավարկություն։

\textbf{Կեղծ համաձայնության դեմ պաշտպանություն.} Փակ խմբերի միջև փոխադարձ հաստատման օրինաչափությունները հայտնաբերելի են սոցիալական գրաֆի վերլուծության միջոցով։ Հեղինակության ագրեգացման ֆունկցիաները զեղչում են սերտորեն կլաստերացված աղբյուրներից եկող պնդումները։

\textbf{Պետական մակարդակի հակառակորդ.} Սոխաձև երթուղավորումը, կենտրոնացված ենթակառուցվածքի բացակայությունը և Ստուգողի դերի բաշխումը մասնակիցների բազայի վրա ապահովում են, որ ճնշումը պահանջում է միջոցներ, որոնք անհամաչափ են ցանկացած օգուտի։

\textbf{Գաղտնիություն ընդդեմ հաշվետվողականության.} Կեղծանվանականությունը՝ միջին ուղի անանունության և ինքնության բացահայտման միջև, թույլ է տալիս DID-ներին կուտակել իմաստալից հեղինակություն՝ առանց բացահայտելու Կրողի իրական աշխարհի ինքնությունը։

% ============================================================
% 11. PRIVACY
% ============================================================
\section{Գաղտնիության նկատառումներ}

RSN-ի գաղտնիության մոդելը կառուցված է կեղծանվանականության վրա։ Կրողը վերահսկում է ինքնության բացահայտումը՝ նվազագույն (զուտ կեղծանուն), ընտրողական (կապված կոնկրետ հավատարմագրեր), կամ լիարժեք (կապված իրավական ինքնություն)։ Հրապարակված պնդումները մշտական են. համակարգն աջակցում է համատեքստային ուղղմանը, բայց ոչ ջնջմանը։ Կրողը կարող է լքել վտանգված DID-ը և սկսել նորից՝ զոհաբերելով կուտակված հեղինակությունը։

% ============================================================
% 12. RELATED WORK
% ============================================================
\section{Հարակից աշխատանքներ}

W3C DID Core բնութագիրը~\cite{did2022} և Ceramic Network-ը~\cite{ceramic2023} տրամադրում են ինքնության հիմքը։ EigenTrust-ը~\cite{kamvar2003} ցուցադրեց բաշխված հեղինակության ագրեգացումը՝ առանց կենտրոնական իշխանության։ SBT-ները~\cite{weyl2022} ներմուծեցին չփոխանցելի սոցիալական նշաններ։ RSN-ը տարբերվում է տնտեսական ծախսի՝ որպես որակի զտիչի շեշտադրմամբ և արմատականությունը գնագոյացնելու իր մեխանիզմով։

DAO-ները~\cite{buterin2014} և քառակուսի քվեարկությունը~\cite{buterin2019} ցուցադրեցին ապակենտրոնացված կառավարման իրագործելիությունը։ RSN-ը համաձայնությունը բխեցնում է երկկողմ գնային բանակցություններից, ոչ թե քվեարկությունից։

Առաջարկը հենվում է անարխո-կապիտալիստական տեսության վրա~\cite{rothbard1973,hoppe2001}՝ մասնավոր ծառայությունների մատուցման համար, բայց շեղվում է երկու կրիտիկական առումով. այն նախագծում է չարության և վրիժառու մղումների համար՝ ոչ թե ենթադրելով բանականություն, և առաջարկում է աստիճանական անցում՝ ոչ թե հեղափոխություն։ Առաջացող հասարակական պայմանագրերի հասկացությունն ունի նախադեպեր Հայեկի ինքնաբուխ կարգի տեսության մեջ~\cite{hayek1973}։

\textbf{Անարխո-ագորիզմ.} RSN-ը կիսում է զգալի ընդհանուր հիմք ագորիստական հակատնտեսագիտության հետ~\cite{konkin2006}՝ հատկապես զուգահեռ ինստիտուտների կառուցման և կամավոր փոխանակման շեշտադրումով։ Այնուամենայնիվ, ագորիզմը հակված է ենթադրելու բանական անձնական շահը որպես բավարար համակարգման մեխանիզմ և հաճախ չունի կոնկրետ միգրացիոն ուղի գոյություն ունեցող պետական կառույցներից։ RSN-ը հստակորեն ներառում է ոչ բանական վարքագծային հակումները և ապահովում ֆիսկալ ճնշման մեխանիզմ աստիճանական անցման համար։

\textbf{Մերիտոկրատիա.} RSN-ը կարող է ներկայացնել առաջին գործառական նկարագրությունն այն բանի, թե ինչպես կարող է մերիտոկրատիան~\cite{young1958} առաջանալ որպես համակարգային հատկություն, ոչ թե կարգախոս։ Ավանդական մերիտոկրատական համակարգերը ձախողվում են, քանի որ պահանջում են կենտրոնական իշխանություն «արժանիքը» սահմանելու և կիրառելու համար։ RSN-ում արժանիքն առաջանում է երկկողմ գնահատումներից. նրանք, ովքեր արժեք են ստեղծում, կուտակում են հեղինակություն. ոչ մի հանձնաժողով չի որոշում, թե ով ունի արժանիք։

\textbf{Սեփականությունը որպես արտոնություն.} RSN-ը հիմնարար կերպով շեղվում է սեփականության իրավունքների անարխո-կապիտալիստական և ավստրիական-դպրոցի վերաբերմունքից՝ որպես բնական և բացարձակ։ RSN շրջանակում սեփականությունը \emph{արտոնություն} է՝ սոցիալական ճանաչում, որը կարող է, ծայրահեղ դեպքերում, հետկանչվել համայնքի կողմից, երբ մասնակիցը հիմնարար կերպով խախտում է ընդհանուր նորմերը (դավաճանություն, գոյաբանական ճգնաժամում համայնքը պաշտպանելուց հրաժարվել, համակարգված շահագործում)։ Սա պետական բռնագրավում չէ, այլ սոցիալական ճանաչման հետկանչ երկկողմ հարաբերությունների ցանցի կողմից։

% ============================================================
% 13. LIMITATIONS
% ============================================================
\section{Քննարկում և սահմանափակումներ}

\textbf{Սառը մեկնարկ.} RSN-ի արժեքը կախված է ցանցային էֆեկտներից. վաղ ընդունողները բախվում են սահմանափակ արժեքի, մինչև մասնակցությունը հասնում է շեմի։ Այնուամենայնիվ, նույնիսկ վաղ փուլերից DID ցանցը ծառայում է որպես օգտակար լրացուցիչ տեղեկատվության աղբյուր։ Ակնկալվում է, որ վաղ հեղինակության գնահատումը և իշխանության գնահատումը կզարգանան աստիճանաբար՝ աճող բարդությամբ։

\textbf{Հակամակաբուծություն և հասանելիության վերահսկում.} Թիրախ DID-ները կարող են աստիճանական վճարներ և հասանելիության սահմանափակումներ պարտադրել ցածր հեղինակություն ունեցող DID-ների հարցումների վրա։ Սա բինար չէ, այլ շարունակական սանդղակ. որքան քիչ հեղինակություն ունի հարցում ներկայացնող DID-ը, այնքան քիչ տեղեկատվություն է ստանում, այնքան երկար է սպասում և այնքան ավելի է վճարում։ Այս մեխանիզմը խրախուսում է իսկական մասնակցությունը պասիվ սպառման փոխարեն։

\textbf{Հասանելիության անհավասարություն.} Պնդումներ հրապարակելու արժեքը կարող է զտել օրինական բողոքները տնտեսապես անապահով մասնակիցներից։ Մեղմացում. յուրաքանչյուր մասնակից միաժամանակ ծառայում է որպես պոտենցիալ ստուգող (կամ պատվիրակում է այս դերը) և վաստակում ստուգման վճարներ։ Բավական երկար ժամանակահատվածում ոչ արմատական մասնակիցը պետք է մոտենա ֆինանսական չեզոքության՝ ստուգման եկամուտը մոտավորապես հավասարակշռելով հրապարակման ծախսերին։ Սա վիճակագրական ակնկալիք է, ոչ երաշխիք։

\textbf{Հեղինակության անհավասարություն.} Վաղ ընդունողները կուտակում են առավելություններ, որոնք կարող են խոչընդոտներ ստեղծել ուշ եկողների համար։ Այնուամենայնիվ, հեղինակության գնահատումն իրականացվում է երրորդ կողմի մատակարարների կողմից, ովքեր կարող են ընտրել մեղմել առաջին քայլողի առավելությունն իրենց ագրեգացման ալգորիթմների միջոցով։ Լավ հեղինակությամբ առաջինը լինելը օրինական համեմատական տնտեսական առավելություն է, և դա նախագծմամբ է։

\textbf{Բաց հարցեր.} Օպտիմալ ծախսային ֆունկցիաները, ագրեգացման ստանդարտները, արձանագրության կառավարումը և AI-ի կողմից գեներացված սինթետիկ հեղինակության տվյալների հետ փոխազդեցությունը մնում են ապագա աշխատանքի ոլորտներ։

Այս հոդվածը չի պնդում, որ RSN-ը կարտադրի օպտիմալ արդյունքներ բոլոր հանգամանքներում։ Այն պնդում է, որ RSN-ը ներկայացնում է \emph{իրագործելի} այլընտրանք՝ տեխնիկապես իրականացվող, տնտեսապես ինքնակայուն և սոցիալապես համատեղելի մարդկային վարքագծային հակումների հետ այնպես, ինչպես դրանք իրականում կան։

% ============================================================
% 14. CONCLUSION
% ============================================================
\section{Եզրակացություն}

Այս հոդվածը ներկայացրեց Հեղինակության սոցիալական ցանցը՝ ապակենտրոնացված արձանագրություն, որը հնարավոր է դարձնում կեղծանվանական, ստուգելի հեղինակության փոխանակումը։ Թանկ արմատականության սկզբունքն ապահովում է, որ խարդախ պնդումները գնագոյացվեն դուրս՝ առանց գրաքննության։ Առաջարկվող միգրացիոն ուղին հնարավոր է դարձնում աստիճանական, հետշրջելի անցումը գոյություն ունեցող ինստիտուտներից։ Նախագիծը ներառում է մարդկային բնությունն այնպես, ինչպես կա՝ ներառյալ մանրախնդրությունը, չարությունը և վրիժառու բավարարվածության ցանկությունը, փոխանակ պահանջելու գաղափարախոսական վերափոխում։ Արդյունքը ուտոպիա չէ, այլ համակարգ, որտեղ արդարությունը հասանելի է, հեղինակությունը վաստակվում է, և իրավախախտման արժեքը կրում է այն կողմը, որը պատճառել է այն։

% ============================================================
% REFERENCES
% ============================================================
\bibliography{references}

\end{document}
